\begin{codebox}
\codeline{\textcolor{cmtgray}{\#\ AI时代的本体定位：幻觉控制与知识权威源}}
\codeline{\textcolor{cmtgray}{\#\ Ontology\ in\ the\ AI\ Era:\ Hallucination\ Control\ and\ Source\ of\ Truth}}
\codeline{\textcolor{cmtgray}{\#}}
\codeline{\textcolor{cmtgray}{\#\ 大语言模型（LLM）改变了知识获取方式，也带来了新问题：幻觉}}
\codeline{\textcolor{cmtgray}{\#\ 本体在\ AI\ 时代的核心角色：作为企业知识的"唯一权威来源"}}
\codeblank{}
\codeline{\textcolor{cmtgray}{\#\ ==============================}}
\codeline{\textcolor{cmtgray}{\#\ 1.\ 大模型的知识困境}}
\codeline{\textcolor{cmtgray}{\#\ ==============================}}
\codeline{\textcolor{cmtgray}{\#\ LLM\ 的本质是统计推断：根据上下文预测最可能的下一个词}}
\codeline{\textcolor{cmtgray}{\#\ \ \ 它"知道"的是语言模式，而不是经过核实的事实}}
\codeline{\textcolor{cmtgray}{\#\ \ \ 当训练数据缺失或模糊时，它会生成"看起来合理"的内容——幻觉}}
\codeblank{}
\codeline{\textcolor{cmtgray}{\#\ 幻觉在工程场景中的代价：}}
\codeline{\textcolor{cmtgray}{\#\ \ \ 设备参数答错\ \(\rightarrow\)\ 工艺事故}}
\codeline{\textcolor{cmtgray}{\#\ \ \ 法规条款编造\ \(\rightarrow\)\ 合规风险}}
\codeline{\textcolor{cmtgray}{\#\ \ \ 维修步骤臆造\ \(\rightarrow\)\ 安全事故}}
\codeblank{}
\codeline{\textcolor{cmtgray}{\#\ ==============================}}
\codeline{\textcolor{cmtgray}{\#\ 2.\ 两种回答路径的对照}}
\codeline{\textcolor{cmtgray}{\#\ ==============================}}
\codeline{\textcolor{cmtgray}{\#\ 问题："Lathe\_003\ 的最大功率是多少？可以加工钛合金吗？"}}
\codeblank{}
\codeline{\textcolor{cmtgray}{\#\ 路径A\ ——\ LLM\ 直接回答（无本体）：}}
\codeline{\textcolor{cmtgray}{\#\ \ \ 模型基于训练语料中"车床通常……"的统计模式作答}}
\codeline{\textcolor{cmtgray}{\#\ \ \ \(\rightarrow\)\ 数值可能凭空生成，能力判断无依据，错误无法被发现}}
\codeblank{}
\codeline{\textcolor{cmtgray}{\#\ 路径B\ ——\ 本体约束下回答（有本体）：}}
\codeline{\textcolor{cmtgray}{\#\ \ \ Step\ 1\ \ 实体链接：'Lathe\_003'\ \(\rightarrow\)\ mfg:Lathe\_003（本体中的个体）}}
\codeline{\textcolor{cmtgray}{\#\ \ \ Step\ 2\ \ SPARQL查询事实：}}
\codeline{\textcolor{cmtgray}{\#\ \ \ \ \ \ \ \ \ \ \ SELECT\ ?p\ WHERE\ \{\ mfg:Lathe\_003\ mfg:hasPower\ ?p\ \}}}
\codeline{\textcolor{cmtgray}{\#\ \ \ \ \ \ \ \ \ \ \ \(\rightarrow\)\ 15.5\ kW（来自设备台账，有出处）}}
\codeline{\textcolor{cmtgray}{\#\ \ \ Step\ 3\ \ 能力校验：}}
\codeline{\textcolor{cmtgray}{\#\ \ \ \ \ \ \ \ \ \ \ ASK\ \{\ mfg:Lathe\_003\ mfg:canProcess\ mfg:TitaniumAlloy\ \}}}
\codeline{\textcolor{cmtgray}{\#\ \ \ \ \ \ \ \ \ \ \ \(\rightarrow\)\ false（本体中无此断言）}}
\codeline{\textcolor{cmtgray}{\#\ \ \ Step\ 4\ \ 生成回答：功率\ 15.5\ kW（来源：设备台账\ 2026-03）；}}
\codeline{\textcolor{cmtgray}{\#\ \ \ \ \ \ \ \ \ \ \ 当前知识库中没有该设备可加工钛合金的记录（区别于"不能"）}}
\codeblank{}
\codeline{\textcolor{cmtgray}{\#\ 关键差异：}}
\codeline{\textcolor{cmtgray}{\#\ \ \ 路径B的每个事实都可溯源到结构化知识；}}
\codeline{\textcolor{cmtgray}{\#\ \ \ "知识库中没有记录"与"事实为否"被严格区分（开放世界假设，见第2章）}}
\codeblank{}
\codeline{\textcolor{cmtgray}{\#\ ==============================}}
\codeline{\textcolor{cmtgray}{\#\ 3.\ 本体作为幻觉控制层}}
\codeline{\textcolor{cmtgray}{\#\ ==============================}}
\codeline{\textcolor{cmtgray}{\#\ 本体为\ AI\ 系统提供三类硬约束：}}
\codeblank{}
\codeline{\textcolor{cmtgray}{\#\ (1)\ 词汇约束\ ——\ AI\ 只能使用本体中定义的类与属性}}
\codeline{\textcolor{cmtgray}{\#\ \ \ \ \ 防止编造不存在的概念（"超导车床"）}}
\codeblank{}
\codeline{\textcolor{cmtgray}{\#\ (2)\ 公理约束\ ——\ 违反公理的陈述直接被拒绝}}
\codeline{Running\ \(\sqcap\)\ Maintenance\ \(\sqsubseteq\)\ \(\bot\)\ \ \ \ \ \ \ \ \#\ "设备正在运行中维护"是非法陈述}
\codeline{Equipment\ \(\sqcap\)\ Material\ \(\sqsubseteq\)\ \(\bot\)\ \ \ \ \ \ \ \ \ \ \#\ "把车床作为原材料投料"是非法陈述}
\codeblank{}
\codeline{\textcolor{cmtgray}{\#\ (3)\ 范围约束\ ——\ 数值必须落在物理合理区间}}
\codeline{\textcolor{cmtgray}{\#\ \ \ \ \ 功率\ ∈\ [0,\ 10000]\ kW；温度\ \(\ge\)\ -273.15\ °C}}
\codeline{\textcolor{cmtgray}{\#\ \ \ \ \ （SHACL\ 实现见第7章质量校验）}}
\codeblank{}
\codeline{\textcolor{cmtgray}{\#\ ==============================}}
\codeline{\textcolor{cmtgray}{\#\ 4.\ LLM\ 与本体的互补结构}}
\codeline{\textcolor{cmtgray}{\#\ ==============================}}
\codeline{\textcolor{cmtgray}{\#\ \ \ LLM\ 擅长：自然语言理解、意图识别、文本生成、模糊匹配}}
\codeline{\textcolor{cmtgray}{\#\ \ \ 本体擅长：事实存储、一致性检验、逻辑推理、可解释溯源}}
\codeblank{}
\codeline{\textcolor{cmtgray}{\#\ \ \ 分工模式：}}
\codeline{\textcolor{cmtgray}{\#\ \ \ 用户自然语言\ \(\rightarrow\)\ [LLM]\ 解析意图、生成查询}}
\codeline{\textcolor{cmtgray}{\#\ \ \ \ \ \ \ \ \ \ \ \ \ \ \ \(\rightarrow\)\ [本体/知识图谱]\ 检索事实、校验约束、执行推理}}
\codeline{\textcolor{cmtgray}{\#\ \ \ \ \ \ \ \ \ \ \ \ \ \ \ \(\rightarrow\)\ [LLM]\ 把结构化结果组织成自然语言回答}}
\codeline{\textcolor{cmtgray}{\#\ \ \ （实现模式详见第8章：GraphRAG、Text2SPARQL、本体引导Agent）}}
\codeblank{}
\codeline{\textcolor{cmtgray}{\#\ ==============================}}
\codeline{\textcolor{cmtgray}{\#\ 5.\ 工业界实践方向}}
\codeline{\textcolor{cmtgray}{\#\ ==============================}}
\codeline{\textcolor{cmtgray}{\#\ 数据分析平台：以本体层（Ontology\ Layer）作为企业数据的语义中枢，}}
\codeline{\textcolor{cmtgray}{\#\ \ \ 业务对象、关系、行动在本体中统一定义，AI\ 在其上运行}}
\codeline{\textcolor{cmtgray}{\#\ 制造企业：设备、工艺、质量知识图谱，支撑故障诊断与工艺推荐}}
\codeline{\textcolor{cmtgray}{\#\ 监管行业：法规条款本体化，让合规判断可溯源、可审计}}
\codeblank{}
\codeline{\textcolor{cmtgray}{\#\ 共同点：让\ AI\ 的每个结论都能追溯到一条结构化的知识定义，}}
\codeline{\textcolor{cmtgray}{\#\ 从根源上抑制大模型幻觉——这正是本书后续各章要构建的能力}}
\end{codebox}
